\subsection{Vector fields}

Let X be a manifold of class \(C^{r}\) , and let U be an open set of X.

8.2.1. A section (not necessarily continuous, cf. Notations and Conventions) of the tangent bundle T(X) over U is called a vector field on U.

8.2.2. The dual bundle T'(X) (7.7.3) of T(X) is called the cotangent bundle of X. Its sections are called covector fields. If f is a function of class \(C^{k}\) on X (with \(1 \leqslant k \leqslant r\) ), with values in K, its differential df: \(x \mapsto d_{x}f\) (5.5.6) is a covector field of class \(C^{k-1}\) . More generally, if f is a function of class \(C^{k}\) on X ( \(1 \leqslant k \leqslant r\) ), with values in a Banach space E, its differential df is a section of class \(C^{k-1}\) of the bundle \(\mathcal{L}(\mathrm{T}(\mathrm{X}); \mathrm{E}_{\mathrm{X}})\) .

8.2.3. Let \(\xi\) be a vector field of class \(\mathbf{C}^{r-1}\) on X, and let \(f\in\mathcal{C}^{r}(X;F)\) be a function of class \(\mathbf{C}^{r}\) with values in a Banach space F. We denote by \(\langle\xi,df\rangle\) , or \(\mathbf{D}_{\xi}(f)\) , or again \(\xi(f)\) the function \(x\mapsto d_{x}f(\xi(x))\) ; it is an element of \(\mathcal{C}^{r-1}(X;F)\) . For \(f\) fixed, the map \(\xi\mapsto\mathbf{D}_{\xi}(f)\) is \(\mathcal{C}^{r-1}(X)\) -linear. If \(\mathbf{D}_{\xi}f=0\) for every function \(f\) of class \(\mathbf{C}^{r}\) with values in a Banach space, defined on an open set of X, we have \(\xi=0\) .

Let E, F, G be Banach spaces and let \((u, v) \mapsto u \cdot v\) be a continuous bilinear map from \(\mathbf{E} \times \mathbf{F}\) to \(\mathbf{G}\) . Let \(f \in \mathcal{C}^r(\mathbf{X}; \mathbf{E})\) , \(g \in \mathcal{C}^r(\mathbf{X}; \mathbf{F})\) and let \(f.g \in \mathcal{C}^r(\mathbf{X}; \mathbf{G})\) be their product. If \(\xi\) is a vector field of class \(\mathbf{C}^{r-1}\) on \(\mathbf{X}\) , one has

\[
\mathrm{D} _ {\xi} (f. g) = \mathrm{D} _ {\xi} (f). g + f. \mathrm{D} _ {\xi} (g).
\]

In particular, the map \(\mathbf{D}_{\xi}\colon\mathcal{C}^{r}(\mathbf{X})\to\mathcal{C}^{r-1}(\mathbf{X})\) is a derivation of the algebra \(\mathcal{C}^{r}(\mathbf{X})\) into the \(\mathcal{C}^{r}(\mathbf{X})\)-module \(\mathcal{C}^{r-1}(\mathbf{X})\) .

8.2.4. Suppose that X is locally of finite dimension and that one of the following two hypotheses holds:

\begin{enumerate}
\def\labelenumi{(\roman{enumi})}
\item
  \(r = \infty\) and X is Hausdorff;
\item
  \(r = \omega\) and X is isomorphic to an open polydisc of \(K^{n}\) .
\end{enumerate}

Then, for every derivation D of the algebra \(\mathcal{C}^r(X)\) , there exists one and only one vector field \(\xi\) of class \(C^r\) such that \(D = D_{\xi}\) .

8.2.5. Let \((u^{1},\ldots ,u^{n})\) be a coordinate system of X on U, and \(\xi\) a vector field on U. The \(i\)-th coordinate of \(\xi\) in the tangent frame (8.1.5) defined by \((u^{1},\ldots ,u^{n})\) is \(\mathrm{D}_{\xi}(u^i)\) . In other words, we have

\[
\xi = \sum_ {1 \leqslant i \leqslant n} \mathrm{D} _ {\xi} (u ^ {i}). \partial / \partial u ^ {i}.
\]

8.2.6. Let \(\varphi: X \to Y\) be a morphism of manifolds of class \(C^r\) , \(\xi\) a vector field on \(X\) and \(\eta\) a vector field on \(Y\) . We say that \(\xi\) is \(\varphi\)-related to \(\eta\) if, for all \(x \in X\) , one has \(\eta(\varphi(x)) = T_x(\varphi)(\xi(x))\) . Suppose that this is so, and that \(\xi\) and \(\eta\) are of class \(C^{r-1}\) . Then, for every Banach space E, the diagram

\[
\begin{array}{c c c} \mathcal {C} ^ {r} (\mathrm{Y}; \mathrm{E}) & \xrightarrow {\varphi^ {*}} & \mathcal {C} ^ {r} (\mathrm{X}; \mathrm{E}) \\ \Big \downarrow_ {\mathrm{D} _ {\eta}} & & \Big \downarrow_ {\mathrm{D} _ {\xi}} \\ \mathcal {C} ^ {r - 1} (\mathrm{Y}; \mathrm{E}) & \xrightarrow {\varphi^ {*}} & \mathcal {C} ^ {r - 1} (\mathrm{X}; \mathrm{E}) \end{array}
\]

\(\mathrm{where}\ \varphi^{*}(f) = f\circ \varphi ,\) is commutative.

When \(\varphi\) is étale, for every vector field \(\eta\) on Y, there exists one and only one vector field on X that is \(\varphi\) -related to \(\eta\) ; we denote it \(\varphi^{*}\eta\) .

8.2.7. Let \(\varphi: X \to Y\) be a morphism of manifolds of class \(C^r\) . If \(\omega\) is a field of covectors \(^{(1)}\) on \(Y\), one defines a field of covectors \(\varphi^*\omega\) on \(X\) by

\[
(\varphi^ {*} \omega) (x) = ^ {t} (\mathrm{T} _ {x} (\varphi)) (\omega (\varphi (x))) (x \in X).
\]

More generally, let \(\tau\) be a vector functor (7.6) of class \(\mathbf{C}^{r-1}\) , of type \((\mathbf{I}_{+}, \mathbf{I}_{-})\) with \(\mathbf{I}_{+} = \varnothing\) and \(\mathbf{I}_{-}\) reduced to a single element (such a functor is said to be contravariant). If \(\omega\) is a section of class \(\mathbf{C}^{r-1}\) of the vector bundle \(\tau(\mathrm{T}(\mathrm{Y}))\) with base \(\mathbf{Y}\) we define a section \(\varphi^{*}\omega\) of \(\tau(\mathrm{T}(\mathrm{X}))\) by \((\varphi^{*}\omega)(x) = \tau(\mathrm{T}_{x}(\varphi))(\omega(\varphi(x)))\) ( \(x \in \mathrm{X}\) ).
% semantic-fragments: legacy-input-seam
\relax{}
